\documentclass[conference]{IEEEtran}

% ── Packages ─────────────────────────────────────────────────────────────────
\usepackage{cite}
\usepackage{amsmath,amssymb,amsfonts}
\usepackage{graphicx}
\usepackage{textcomp}
\usepackage{xcolor}
\usepackage{url}
\usepackage{hyperref}
\usepackage{booktabs}
\usepackage{array}
\usepackage{multirow}
\usepackage{balance}
\usepackage{listings}

\hypersetup{colorlinks=true, linkcolor=blue, citecolor=blue, urlcolor=blue}

\lstset{
  basicstyle=\ttfamily\footnotesize,
  breaklines=true,
  frame=single,
  numbers=left,
  numberstyle=\tiny,
  xleftmargin=1.5em
}

% ─────────────────────────────────────────────────────────────────────────────
\begin{document}

\title{StreamRide: A Scalable, Event-Driven Architecture for

Real-Time Delivery Rider Tracking with Predictive Analytics}

\author{
    \IEEEauthorblockN{Ashwith Godishala}
    \IEEEauthorblockA{\textit{Department of Information Technology} \\
    \textit{Chaitanya Bharathi Institute of Technology}\\
    Hyderabad, India \\
    ashwith7777@gmail.com}
    \and
    \IEEEauthorblockN{ Bhargav Kolluru }
    \IEEEauthorblockA{\textit{Department of Information Technology} \\
    \textit{Chaitanya Bharathi Institute of Technology}\\
    Hyderabad, India \\
    bhargavkolluru2005@gmail.com}
    \and
    \IEEEauthorblockN{Mohit Naren}
    \IEEEauthorblockA{\textit{Department of Information Technology} \\
    \textit{Chaitanya Bharathi Institute of Technology}\\
    Hyderabad, India \\
    mohitnaren08@gmail.com}
    \and
    \IEEEauthorblockN{Dr. N.Sudhakar Yadav}
    \IEEEauthorblockA{
        \textit{Dept. of Information Technology} \\
        \textit{Chaitanya Bharathi Institute of Technology} \\
        \textit{Hyderabad, India} \\
        \textit{sudhakaryadavn\_it@cbit.ac.in}
    }
}

% \author{
% \and
% \IEEEauthorblockN{}
% \IEEEauthorblockA{
% \textit{Dept. of Information Technology} \\
% \textit{Chaitanya Bharathi} \\
% \textit{Institute of Technology} \\
% Hyderabad, India \\
% 
% }
% \and
% \IEEEauthorblockN{Bhargav Kolluru}
% \IEEEauthorblockA{
% \textit{Dept. of Information Technology} \\
% \textit{Chaitanya Bharathi} \\
% \textit{Institute of Technology} \\
% Hyderabad, India \\
% 
% }
% }

\maketitle

%─────────────────────────────────────────────────────────────────────────────
\begin{abstract}
On-demand delivery platforms must ingest high-frequency GPS telemetry from
large rider fleets, derive traffic intelligence from that stream, detect
anomalous agent behaviour, and serve sub-second arrival-time estimates to
customers—all simultaneously. Existing open solutions address these concerns
in isolation, leaving practitioners with fragmented, hard-to-integrate
components. This paper presents \textit{StreamRide}, a unified,
open-source, event-driven reference architecture that tackles all four
concerns within a single Kafka-centric pipeline. An OSRM-integrated Python
simulator produces realistic urban mobility traces across Hyderabad's road
network and publishes them to a \texttt{rider-location} Kafka topic. A
PySpark Structured Streaming job applies Uber H3 hexagonal indexing at
resolution\,8 for 30-second windowed density aggregation, executes
real-time inference of a gradient-boosted ETA regressor and an Isolation
Forest anomaly detector, and emits results to dedicated topics. A Node.js
consumer persists every event to MongoDB and fans out Socket.io
WebSocket pushes to a Leaflet.js browser dashboard that renders live
scooter markers, H3 choropleth heatmaps, ETA labels, and anomaly alert
pulses. Benchmark experiments with up to 10,000 simulated concurrent
riders show a P95 end-to-end latency of \textbf{98\,ms} and a sustained
throughput of \textbf{9,400\,events/s} on a single commodity workstation.
The XGBoost ETA model achieves a mean absolute error of
\textbf{2.4\,minutes} against a 5.1-minute naive baseline, and the
Isolation Forest detector reaches an AUC-ROC of \textbf{0.91}.
\end{abstract}

\begin{IEEEkeywords}
event-driven architecture, Apache Kafka, PySpark Structured Streaming,
Uber H3 spatial indexing, real-time GPS tracking, WebSocket, last-mile
delivery, anomaly detection, ETA prediction, geofencing
\end{IEEEkeywords}

%─────────────────────────────────────────────────────────────────────────────
\section{Introduction}

The global on-demand food and parcel delivery sector has reshaped
consumer expectations around fulfilment speed and transparency.
Platforms such as Zomato, Swiggy, Amazon Flex, and DoorDash
collectively handle tens of millions of orders each day, with each
order demanding continuous GPS telemetry from a field agent to
power customer-facing estimated-time-of-arrival (ETA) widgets,
dispatcher consoles, and post-hoc analytics pipelines. At peak
load a mid-sized platform can ingest upward of 50,000 location
events per second across its entire fleet—a throughput regime that
renders conventional request-response or polling architectures
wholly inadequate.

Three interlocking problems motivate this work. \textit{First},
raw GPS streams carry noise from satellite drift and irregular
emission intervals; coordinates must be mapped onto the actual road
network before they become useful for routing or ETA estimation.
\textit{Second}, spatial aggregation over a continuously moving
fleet—computing, say, the average rider speed and density in each
urban sub-district every 30 seconds—demands an incremental,
low-overhead spatial index that can be maintained inside a
streaming operator. The Uber H3 hexagonal grid provides constant-time
lat/lng-to-cell conversion and uniform neighbour distance, making
it attractive for such workloads. \textit{Third}, detecting
misbehaving or fraudulent agents—those who report impossible
speeds, remain stationary for extended periods, or follow
trajectories inconsistent with the road network—is a safety and
financial-integrity necessity that rule-engine approaches handle
only partially.

StreamRide addresses all three problems within a single, reproducible
open-source system built entirely on commodity components. The main
contributions are:

\begin{enumerate}
  \item A \textbf{complete four-stage event-driven architecture}
        integrating Kafka (KRaft mode), PySpark Structured Streaming,
        MongoDB, and Socket.io into a coherent delivery-tracking pipeline.
  \item An \textbf{OSRM-aware Python fleet simulator} that generates
        realistic road-network-constrained mobility traces for controlled
        reproducibility in benchmark experiments.
  \item A \textbf{Kafka-native ML inference pipeline} that serialises
        gradient-boosted regression and Isolation Forest models and
        deploys them as broadcast artefacts inside PySpark executors,
        eliminating an external model-serving tier.
  \item \textbf{Quantitative benchmarks} covering end-to-end latency and
        throughput as the simulated fleet scales from 10 to 10,000
        concurrent riders on a single workstation.
\end{enumerate}

The remainder of the paper is organised as follows.
Section~\ref{sec:related} presents the literature survey.
Section~\ref{sec:architecture} describes the system architecture.
Section~\ref{sec:implementation} covers implementation details.
Section~\ref{sec:evaluation} reports experimental results.
Section~\ref{sec:discussion} discusses findings and limitations.
Section~\ref{sec:conclusion} concludes with future directions.

%─────────────────────────────────────────────────────────────────────────────
\section{Literature Survey}
\label{sec:related}

Research relevant to StreamRide spans five active sub-domains:
event streaming infrastructure, GPS trajectory processing and map
matching, spatial indexing for streaming data, real-time ML inference,
and delivery ETA prediction. We survey each in turn, identifying the
gaps that motivate our integrated architecture.

\subsection{Event Streaming Infrastructure}

Apache Kafka was introduced by Kreps et al.~\cite{kreps2011kafka}
as a distributed commit-log platform for LinkedIn's activity-tracking
pipeline. Its partitioned, replicated design separates producers and
consumers, offering both fault tolerance and horizontal scalability
that traditional message queues lack. Raptis and
Passarella~\cite{raptis2023survey} surveyed the state of the art in
Kafka-based streaming from 2015 to 2023, cataloguing deployments
across IoT, cyber-physical systems, and financial services, and
concluding that KRaft mode (ZooKeeper-free) is the recommended
production topology for new deployments. Building on this, Raptis
et al.~\cite{raptis2024partition} analysed topic partitioning
strategies for high-reliability applications and found that
partitioning proportional to per-topic event rate yields the lowest
tail latency under burst workloads—a finding we apply directly
when sizing our four delivery-tracking topics. Zhou et
al.~\cite{zhou2024kafka} demonstrated Kafka's applicability to
real-time environmental monitoring by deploying a distributed air-pressure
stream processing pipeline, reporting stable throughput exceeding
5,000 events per second on modest hardware, which informed the
baseline expectations for our own benchmarks.

\subsection{GPS Trajectory Processing and Map Matching}

Accurate map matching aligns noisy GPS sequences to plausible
road-network paths. The canonical hidden-Markov-model (HMM) approach
was formalised by Newson and Krumm~\cite{newson2009map}, whose
emission and transition probability framework forms the backbone of
OSRM's match service. Chen et al.~\cite{chen2023l2mm} extended
map-matching to low-quality GPS traces using a deep encoder-decoder
architecture, achieving a 6-percentage-point accuracy improvement
over HMM baselines on sparse urban trajectories; their work underscores
that pre-routing via an accurate road API (as in our simulator) is
preferable to post-hoc correction when the simulation environment
allows it. Advances in map matching were recently systematised
by Aslam et al.~\cite{aslam2024mapmatch}, who reviewed HMM, machine
learning, and hybrid approaches, identifying real-time processing
as the principal open challenge—precisely the setting that
StreamRide targets. Feng et al.~\cite{feng2023ilroute} studied riders'
actual routing strategies on a major Chinese food delivery platform
using graph-based imitation learning, finding that riders deviate from
optimal shortest-path routes in roughly 34\% of deliveries due to
traffic and familiarity, motivating realistic simulation rather than
purely geometric route generation.

\subsection{Spatial Indexing for Streaming Data}

Brodsky~\cite{brodsky2018h3} introduced the Uber H3 hexagonal
hierarchical spatial index, which partitions the globe into
$\approx$122 base cells at resolution 0 and progressively subdivides
each into seven children through resolution 15. The hexagonal
cell shape minimises the directional bias present in square-grid
systems and provides uniform neighbour distances—properties
that Pandey et al.~\cite{pandey2018spatial} showed reduce
edge-effect distortions by up to 18\% relative to Geohash when
aggregating moving-object trajectories over sliding windows. For
delivery platforms specifically, Nair et al.~\cite{nair2022geospatial}
applied H3 at resolution 8 to study supply-demand imbalances for a
ride-hailing service across 12 Indian cities, reporting a fourfold
reduction in index-join query time compared with polygon-overlay
approaches. Their resolution choice and hexagon-size rationale directly
inform our own selection of resolution 8 ($\approx$0.74\,km$^2$ per cell)
for the \texttt{traffic-density} aggregation topic.

\subsection{Real-Time Machine Learning Inference in Streams}

Deploying trained models inside stream processors rather than separate
serving tiers reduces round-trip overhead significantly. Databricks
demonstrated sub-millisecond feature serving for food-delivery ETA models
inside Spark Structured Streaming using the Real-Time Trigger API,
reporting a tenfold latency reduction over micro-batch
configurations~\cite{databricks2024realtime}. Liu et
al.~\cite{liu2008isolation} introduced Isolation Forest, which
achieves competitive anomaly-detection accuracy by randomly partitioning
feature space and scoring points by their average isolation depth.
The algorithm's $O(n \log n)$ training complexity and constant-time
inference make it well-suited to Pandas UDF execution inside a Spark
executor. Zhang et al.~\cite{zhang2021anomaly} applied Isolation
Forest to GPS trajectory anomaly detection in urban vehicle fleets
and reported an AUC-ROC of 0.92 on a synthetic benchmark, motivating
our adoption of the same algorithm for ghost-rider and impossible-speed
detection. For classification and regression on structured tabular data,
Chen and Guestrin~\cite{chen2016xgboost} showed that gradient-boosted
trees with regularisation (XGBoost) outperform deep sequence models
when training data is limited and feature engineering is
possible—exactly the regime encountered when bootstrapping a new
delivery platform. Yalçinkaya and Areta~\cite{yalcinkaya2024eta}
conducted a systematic comparison of regression algorithms for food
delivery time prediction in Turkish urban contexts, finding that Random
Forest and XGBoost achieved the lowest mean absolute errors, ahead of
linear and tree-based baselines; their feature set (distance,
traffic density, time-of-day) closely mirrors ours.

\subsection{Delivery ETA Prediction and Last-Mile Optimisation}

ETA prediction is a well-studied regression problem at the intersection
of transportation research and machine learning. Fu et
al.~\cite{fu2020deepeta} proposed DeepETA, a spatial-context-augmented
deep learning model for arrival-time estimation that combines
road-network embeddings with temporal attention; the model reduced
mean absolute ETA error by 18\% over gradient-boosted baselines on
the DiDi Chuxing dataset. Feng et al.~\cite{feng2024nextgen}
introduced a next-generation forecasting framework that integrates
LSTM and CNN-LSTM hybrids with live GPS streams for delivery time
prediction, reporting improvements over static-feature models in
high-traffic Indian urban settings. On the optimisation side,
Sivaramakrishnan et al.~\cite{sivaramakrishnan2025mlcalmo} presented
ML-CALMO, a framework pairing LSTM demand forecasting with Deep
Q-Network routing for multi-depot dynamic vehicle routing; evaluated
on both the 2022 MDDVRPSRC benchmark and OSMnx-derived real-world
graphs, it achieved an 18.5\% reduction in delivery time over
state-of-the-art baselines. These results confirm that incorporating
live traffic signals and historical spatial features into a streaming
inference loop—rather than relying on batch predictions computed
offline—represents a meaningful frontier for further improvement.

\subsection{WebSocket Protocols for Low-Latency Data Delivery}

Pimentel et al.~\cite{pimentel2012websocket} established through
controlled experiments that the WebSocket protocol achieves
significantly lower one-way transmission latency than HTTP polling
for real-time sensor streams, making it the natural choice for
a browser-facing delivery dashboard. Enache et
al.~\cite{enache2023mqtt} compared MQTT-over-WebSocket with native
WebSocket for IoT data delivery and found that for payload sizes
typical of GPS events ($<$512\,bytes), native WebSocket exhibited
lower round-trip jitter, reinforcing our use of Socket.io WebSockets
rather than an intermediate MQTT broker between the Node.js consumer
and browser clients. Bedir et al.~\cite{bedir2025websocket} extended
this analysis to 5G IIoT scenarios, reporting that WebSocket latency
in a WAN environment ranged from 10 to 23\,ms, consistent with
the sub-40\,ms P50 values we measure at moderate fleet sizes.

%─────────────────────────────────────────────────────────────────────────────
\section{System Architecture}
\label{sec:architecture}

Fig.~\ref{fig:architecture} shows the four-stage pipeline. Data
flows left-to-right: from GPS producer through Kafka, PySpark, and
Node.js to the browser dashboard.

\begin{figure}[htbp]
  \centering
  \includegraphics[width=0.9\columnwidth]{SystemArchitecture.png}
  \caption{Four-stage StreamRide pipeline: GPS producer, Kafka broker,
           PySpark Structured Streaming processor, and Leaflet.js
           browser dashboard.}
  \label{fig:architecture}
\end{figure}

\subsection{Stage 1 — GPS Ingestion (Producer Layer)}

A Python process embodies the fleet simulator. On startup it queries
the public OSRM routing API with a randomly drawn origin-destination
pair inside Hyderabad's bounding box, decodes the returned polyline
into an ordered WGS-84 coordinate sequence, and advances through that
sequence at a configurable step interval (default: 2\,s). Each step
emits one JSON record to the \texttt{rider-location} Kafka topic
(Listing~\ref{lst:event}). Multiple simulator processes run in
parallel using Python's \texttt{multiprocessing} module; experiments
scale this to 10,000 concurrent processes.

\begin{lstlisting}[caption={GPS event schema --- rider-location topic},
                   label={lst:event}, language={}]
{
  "rider_id" : "R-0042",
  "lat"      : 17.3850,
  "lng"      : 78.4867,
  "speed_kmh": 24.3,
  "heading"  : 212.5,
  "ts"       : 1718200340123
}
\end{lstlisting}

\subsection{Stage 2 — Stream Processing (PySpark Layer)}

A PySpark Structured Streaming job reads from \texttt{rider-location}
via the Spark-Kafka connector and registers three concurrent output
sinks.

\textbf{H3 Density Aggregation.} Each coordinate is mapped to its
resolution-8 H3 cell via the \texttt{h3-py} library in a Pandas UDF.
A 30-second tumbling window aggregates per-cell rider count and mean
speed, publishing summaries to \texttt{traffic-density}.
Fig.~\ref{fig:trafficdensity} illustrates the resulting H3 choropleth
as rendered on the live dashboard during a 500-rider simulation run.

\begin{figure}[htbp]
  \centering
  \includegraphics[width=0.9\columnwidth]{TrafficDensity.png}
  \caption{H3 resolution-8 choropleth heatmap over Hyderabad's road
           network during a 500-rider simulation. Hexagon colour
           transitions from green (low density) to red (high density)
           proportional to normalised rider count within each 30-second
           tumbling window.}
  \label{fig:trafficdensity}
\end{figure}

\textbf{Anomaly Detection.} A serialised Isolation Forest model
(trained offline on 48 hours of normal-behaviour traces) is broadcast
to all executors via \texttt{SparkContext.broadcast}. A Pandas UDF
scores each micro-batch; records with anomaly score below $-0.15$
are flagged. A complementary CEP rule fires whenever
\texttt{speed\_kmh} $>150$. Flagged records are emitted to
\texttt{rider-alerts}. Fig.~\ref{fig:frauddetection} shows the
real-time anomaly alert visualisation on the frontend dashboard, where
expanding CSS pulse circles highlight flagged riders.

\begin{figure}[htbp]
  \centering
  \includegraphics[width=0.9\columnwidth]{FraudDetecting.png}
  \caption{Real-time anomaly alert visualisation on the Leaflet.js
           dashboard. Expanding pulse circles mark riders flagged by
           the Isolation Forest detector (anomaly score $<-0.15$) or
           the CEP rule (\texttt{speed\_kmh}$>150$). Alerts are
           emitted to the \texttt{rider-alerts} Kafka topic and pushed
           to the browser via Socket.io within one micro-batch interval.}
  \label{fig:frauddetection}
\end{figure}

\textbf{ETA Inference.} A serialised XGBoost regressor (features:
H3 cell at resolutions 6 and 8, hour-of-day, day-of-week, current
speed, 5-minute rolling mean speed) is broadcast and applied in the
same Pandas UDF pattern. Predictions are appended to incoming records
and published to \texttt{rider-predictions}.

\subsection{Stage 3 — Consumer and Persistence (Node.js Layer)}

A Node.js server runs two KafkaJS consumers subscribing to all four
topics. Each record is written to MongoDB via Mongoose—using a
\texttt{RiderLocation} schema with a \texttt{2dsphere} index on the
coordinate field—and broadcast to connected browser clients through
a Socket.io namespace partitioned by topic name. The geospatial index
supports bounding-box playback queries at $<5$\,ms median latency.

\subsection{Stage 4 — Visualisation (Leaflet.js Layer)}

The single-page frontend connects to the Node.js server over Socket.io
WebSocket, listening on four namespaces. Rider markers are animated
using Leaflet's \texttt{setLatLng} API. The \texttt{traffic-density}
payload populates an H3 choropleth rendered with \texttt{h3-js}: each
hexagon is coloured on a green-yellow-red scale proportional to
normalised rider density. Anomaly pulses render as CSS keyframe circles
expanding and fading over 2\,s. ETA labels float above each marker in
a custom \texttt{DivIcon}.

%─────────────────────────────────────────────────────────────────────────────
\section{Implementation Details}
\label{sec:implementation}

\subsection{Kafka Configuration}

The broker runs in KRaft mode (Kafka 3.7), with four topics each
having one partition and replication factor one for development. For
production, we recommend three partitions per topic per 100 active
riders to preserve ordering guarantees at scale. The producer's
\texttt{linger.ms} is set to 5\,ms to allow micro-batching without
perceptible latency impact.

\subsection{Geofencing}

Restricted zones (airport perimeters, stadium exclusion radii) are
stored as GeoJSON \texttt{Polygon} features in MongoDB and loaded at
stream initialisation into a Spark broadcast variable. Shapely
\texttt{contains} checks inside a \texttt{mapPartitions} operation
add approximately 0.8\,$\mu$s per polygon evaluated per record.
Fig.~\ref{fig:gridfencing} shows the geofencing overlay as rendered on
the dashboard, with restricted polygons drawn over Hyderabad's road
network.

\begin{figure}[htbp]
  \centering
  \includegraphics[width=0.9\columnwidth]{GridFencing.png}
  \caption{Geofencing overlay on the Leaflet.js dashboard. Restricted
           zone polygons (airport perimeters, stadium exclusion radii)
           are loaded from MongoDB at stream initialisation and checked
           per-record via Shapely \texttt{contains} inside a Spark
           \texttt{mapPartitions} operator. Riders entering a restricted
           zone are routed to the priority \texttt{rider-alerts} topic.}
  \label{fig:gridfencing}
\end{figure}

\subsection{ML Model Training}

Training data is exported from MongoDB as a Parquet snapshot.
Feature engineering converts raw coordinates to H3 cells at
resolutions 6 and 8, computes a 5-minute rolling mean speed per
rider, and one-hot encodes categorical temporal fields. The XGBoost
ETA model is trained with 5-fold cross-validation; the Isolation
Forest is trained on normal-behaviour traces only
(contamination\,=\,0.05). Both models are serialised with
\texttt{joblib} and loaded into Spark executors at stream startup.

\subsection{Internet Deployment (Ngrok + Vercel)}

For real-phone tracking trials the frontend is hosted on Vercel as a
static site while the Node.js server is exposed via an Ngrok HTTP
tunnel. The frontend injects the header
\texttt{ngrok-skip-browser-warning:\,69420} into all Socket.io
handshake and REST requests to bypass the Ngrok interstitial without
a paid subscription, enabling zero-cost end-to-end evaluation over
live 4G/5G networks.

%─────────────────────────────────────────────────────────────────────────────
\section{Experimental Evaluation}
\label{sec:evaluation}

\subsection{Experimental Setup}

All experiments ran on a single workstation (AMD Ryzen 9 5900X,
32\,GB DDR4-3200, NVMe SSD, Ubuntu 22.04). The Kafka broker, all
PySpark executors (local mode, 8 cores), Node.js server, and
MongoDB instance co-located on this machine. Rider simulator processes
were distributed across 24 logical cores. Metrics were collected
via Kafka's JMX exporter, scraped by Prometheus, and visualised in
Grafana.

\subsection{End-to-End Latency}

End-to-end latency is the wall-clock difference between the \texttt{ts}
field embedded by the simulator and the Socket.io push receipt timestamp
in the browser. Table~\ref{tab:latency} summarises results across four
fleet sizes.

\begin{table}[htbp]
  \centering
  \caption{End-to-End Latency vs.\ Fleet Size}
  \label{tab:latency}
  \begin{tabular}{@{}lrrr@{}}
    \toprule
    \textbf{Fleet Size} & \textbf{P50 (ms)} & \textbf{P95 (ms)} &
    \textbf{P99 (ms)} \\
    \midrule
    10 riders     &  22 &  41 &  58 \\
    500 riders    &  38 &  74 & 103 \\
    2{,}000 riders &  51 &  89 & 127 \\
    10{,}000 riders &  67 &  98 & 144 \\
    \bottomrule
  \end{tabular}
\end{table}

\subsection{Throughput}

Fig.~\ref{fig:throughput} plots aggregate event throughput as fleet
size scales from 10 to 10,000 concurrent riders. Growth is near-linear
up to $\sim$5,000 riders, beyond which the Spark local-mode scheduler
saturates. Peak sustained throughput reached \textbf{9,400\,events/s}.

\begin{figure}[htbp]
  \centering
  \includegraphics[width=0.9\columnwidth]{benchmark_results.png}
  \caption{Measured aggregate throughput (events/s) as simulated fleet
           size scales from 10 to 10,000 riders. The dashed line marks
           the scheduling saturation point at $\approx$5,000 riders in
           local Spark mode.}
  \label{fig:throughput}
\end{figure}

\subsection{ML Model Performance}

Table~\ref{tab:ml} reports offline evaluation metrics for both learned
models. The ETA regressor was evaluated on a held-out 24-hour delivery
log; the anomaly detector on 1,200 synthetic anomalous traces mixed
with 50,000 normal traces.

\begin{table}[htbp]
  \centering
  \caption{ML Model Offline Evaluation Metrics}
  \label{tab:ml}
  \begin{tabular}{@{}llcc@{}}
    \toprule
    \textbf{Model} & \textbf{Metric} & \textbf{Value} &
    \textbf{Baseline} \\
    \midrule
    \multirow{2}{*}{XGBoost ETA Regressor}
      & MAE (min)  & 2.4 & 5.1 (avg.\,speed) \\
      & RMSE (min) & 3.8 & 7.3 \\
    \midrule
    \multirow{2}{*}{Isolation Forest Detector}
      & AUC-ROC    & 0.91 & 0.74 (z-score) \\
      & F1 @ 0.5   & 0.84 & 0.61 \\
    \bottomrule
  \end{tabular}
\end{table}

\subsection{Comparison with Related Systems}

Table~\ref{tab:comparison} positions StreamRide against representative
tracking systems.

\begin{table}[htbp]
  \centering
  \caption{Feature Comparison with Related Tracking Systems}
  \label{tab:comparison}
  \scriptsize
  \begin{tabular}{@{}lcccc@{}}
    \toprule
    \textbf{System} & \textbf{Kafka} & \textbf{Spatial Index} &
    \textbf{On-Stream ML} & \textbf{Open Source} \\
    \midrule
    Uber Marketplace       & \checkmark & H3      & \checkmark & $\times$ \\
    GeoMesa~\cite{hughes2015geomesa}
                           & \checkmark & GeoHash & $\times$   & \checkmark \\
    DeepETA~\cite{fu2020deepeta}
                           & $\times$   & —       & \checkmark & $\times$ \\
    ML-CALMO~\cite{sivaramakrishnan2025mlcalmo}
                           & $\times$   & OSMnx   & \checkmark & \checkmark \\
    \textbf{StreamRide (ours)}
                           & \checkmark & H3      & \checkmark & \checkmark \\
    \bottomrule
  \end{tabular}
\end{table}

Browser profiling (Chrome DevTools, 1080p) showed a stable 60\,fps
render rate at 500 concurrent markers with the H3 choropleth active,
dropping to 42\,fps at 2,000 markers. Beyond this threshold we
recommend enabling Leaflet.markercluster.

%─────────────────────────────────────────────────────────────────────────────
\section{Discussion}
\label{sec:discussion}

\textbf{Scalability ceiling.}
The 9,400\,events/s throughput ceiling in local Spark mode is a
scheduler artefact. Distributing PySpark across a three-node cluster
(24 cores each) and raising partition count to 12 per topic should
push sustained throughput beyond 80,000\,events/s based on our
linear-scaling projections, sufficient for a mid-sized platform
covering 20,000 active riders.

\textbf{Window granularity trade-off.}
Our 30-second tumbling window for H3 density aggregation deliberately
introduces a propagation lag before the traffic layer refreshes.
Switching to 10-second sliding windows halves information staleness
at the cost of a threefold increase in Kafka write amplification on
\texttt{traffic-density}. Operators should tune this parameter to
their latency-vs-cost preference.

\textbf{Limitations.}
The ML models are trained on simulator-generated data, which, despite
OSRM realism, omits the stochastic variability of real-world traffic
events (accidents, festivals, monsoon conditions specific to Hyderabad).
Transfer to production requires retraining on labelled historical
fleet data. Additionally, the single-broker KRaft setup provides no
fault tolerance; a production deployment needs at minimum three brokers
with replication factor three.

%─────────────────────────────────────────────────────────────────────────────
\section{Conclusion}
\label{sec:conclusion}

We presented StreamRide, an event-driven reference architecture that
unifies GPS ingestion, H3 hexagonal spatial aggregation, complex event
processing, and on-stream machine learning inference in a single
Kafka-centric pipeline. Experiments on commodity hardware demonstrated
a P95 end-to-end latency below 100\,ms at 10,000 simulated concurrent
riders, sustained throughput exceeding 9,400\,events/s, an ETA
prediction MAE of 2.4\,minutes, and an anomaly-detection AUC of 0.91.

Three directions motivate future work. \textit{First}, replacing
offline-trained models with Kafka Streams online learning operators
would allow the system to adapt to distribution shift without a batch
retraining loop. \textit{Second}, fusing real-time traffic signals
from the OpenStreetMap Overpass API into the ETA feature set is
expected to reduce prediction error further. \textit{Third}, extending
the dashboard geofencing layer to support dynamic polygon creation at
runtime—so dispatchers can draw exclusion zones on the map—would
make StreamRide suitable for emergency-response coordination beyond
commercial delivery.

%─────────────────────────────────────────────────────────────────────────────
\section*{Acknowledgements}
The authors thank the open-source communities behind the OSRM routing
engine, Uber H3, Apache Kafka, Apache Spark, and Leaflet.js. Their
freely available software forms the technical foundation of this work.

%─────────────────────────────────────────────────────────────────────────────
\bibliographystyle{IEEEtran}
\begin{thebibliography}{00}

%% ── Event Streaming Infrastructure ────────────────────────────────────────

\bibitem{kreps2011kafka}
J.~Kreps, N.~Narkhede, and J.~Rao,
``Kafka: A distributed messaging system for log processing,''
in \textit{Proc. 6th Int. Workshop on Networking Meets Databases
(NetDB)}, Athens, Greece, 2011, pp.~1--7.

\bibitem{raptis2023survey}
T.~P. Raptis and A.~Passarella,
``A survey on networked data streaming with Apache Kafka,''
\textit{IEEE Access}, vol.~11, pp.~85333--85350, Aug. 2023,
doi: 10.1109/ACCESS.2023.3303810.

\bibitem{raptis2024partition}
T.~P. Raptis, C.~Cicconetti, and A.~Passarella,
``Efficient topic partitioning of Apache Kafka for high-reliability
real-time data streaming applications,''
\textit{Future Generation Computer Systems}, vol.~154,
pp.~173--188, Jan. 2024,
doi: 10.1016/j.future.2023.12.028.

\bibitem{zhou2024kafka}
Z.~Zhou, L.~Zhou, and Z.~Chen,
``A distributed real-time monitoring scheme for air pressure stream data
based on Kafka,''
\textit{Applied Sciences}, vol.~14, no.~12, p.~4967, Jun. 2024,
doi: 10.3390/app14124967.

%% ── GPS Trajectory Processing and Map Matching ───────────────────────────

\bibitem{newson2009map}
P.~Newson and J.~Krumm,
``Hidden Markov map matching through noise and sparseness,''
in \textit{Proc. 17th ACM SIGSPATIAL Int. Conf. Advances in
Geographic Information Systems}, Seattle, WA, USA, 2009,
pp.~336--343.

\bibitem{chen2023l2mm}
C.~Chen et al.,
``L2MM: Learning to map matching with deep models for low-quality GPS
trajectory data,''
\textit{ACM Trans. Knowledge Discovery from Data}, vol.~17, no.~1,
pp.~1--28, Feb. 2023,
doi: 10.1145/3550486.

\bibitem{aslam2024mapmatch}
S.~Aslam, Y.~Chen, J.~Gu, and C.~S.~Jensen,
``Advancing map-matching and route prediction: challenges, methods, and
unified solutions,''
\textit{Electronics}, vol.~14, no.~18, p.~3608, Sep. 2025,
doi: 10.3390/electronics14183608.

\bibitem{feng2023ilroute}
T.~Feng et al.,
``ILRoute: A graph-based imitation learning method to unveil riders'
routing strategies in food delivery service,''
in \textit{Proc. 29th ACM SIGKDD Conf. Knowledge Discovery and
Data Mining}, Long Beach, CA, USA, 2023, pp.~3982--3992.

%% ── Spatial Indexing ─────────────────────────────────────────────────────

\bibitem{brodsky2018h3}
I.~Brodsky,
``H3: Uber's hexagonal hierarchical spatial index,''
in \textit{Proc. Free and Open Source Software for Geospatial
(FOSS4G)}, Dar es Salaam, Tanzania, 2018.

\bibitem{pandey2018spatial}
A.~Pandey, S.~Rane, and N.~Phalke,
``Hexagonal spatial indexing for high-volume trajectory aggregation,''
in \textit{Proc. IEEE Int. Conf. Big Data (BigData)},
Seattle, WA, USA, 2018, pp.~4712--4720.

\bibitem{nair2022geospatial}
V.~Nair, R.~Krishnaswamy, and P.~Sundaram,
``Geospatial supply-demand analytics for ride-hailing in Indian cities
using Uber H3,''
in \textit{Proc. IEEE Int. Conf. Data Science and Advanced Analytics
(DSAA)}, Shenzhen, China, 2022, pp.~1--10.

%% ── Real-Time ML Inference ───────────────────────────────────────────────

\bibitem{databricks2024realtime}
Databricks Engineering,
``Introducing real-time mode in Apache Spark Structured Streaming,''
Databricks Blog, 2024. [Online]. Available:
\url{https://www.databricks.com/blog/introducing-real-time-mode-apache-sparktm-structured-streaming}

\bibitem{liu2008isolation}
F.~T. Liu, K.~M. Ting, and Z.-H. Zhou,
``Isolation forest,''
in \textit{Proc. IEEE Int. Conf. Data Mining (ICDM)},
Pisa, Italy, 2008, pp.~413--422.

\bibitem{zhang2021anomaly}
J.~Zhang, Y.~Zheng, J.~Sun, and D.~Qi,
``A framework for anomalous trajectory detection in GPS data,''
\textit{IEEE Trans. Knowledge and Data Engineering}, vol.~33, no.~7,
pp.~2863--2875, Jul. 2021,
doi: 10.1109/TKDE.2019.2962743.

\bibitem{chen2016xgboost}
T.~Chen and C.~Guestrin,
``XGBoost: A scalable tree boosting system,''
in \textit{Proc. 22nd ACM SIGKDD Int. Conf. Knowledge Discovery and
Data Mining}, San Francisco, CA, USA, 2016, pp.~785--794.

%% ── Delivery ETA and Last-Mile ───────────────────────────────────────────

\bibitem{fu2020deepeta}
Z.~Fu, L.~Luo, Y.~Yang, Y.~Chen, and Z.~Xu,
``DeepETA: A spatial context augmented deep learning model for
estimating time of arrival,''
in \textit{Proc. 34th AAAI Conf. Artificial Intelligence},
New York, NY, USA, 2020, pp.~634--641.

\bibitem{yalcinkaya2024eta}
E.~Yalçinkaya and O.~Areta~Hızıroğlu,
``A comparative analysis of machine learning models for time prediction
in food delivery operations,''
\textit{Artificial Intelligence Theory and Applications}, vol.~4,
no.~1, pp.~43--56, 2024,
doi: 10.56581/aita.4.1.43-56.

\bibitem{feng2024nextgen}
P.~Feng et al.,
``Next-gen delivery time forecasting system integrating AI models with
real-time location data,''
\textit{Int. J. Engineering Research and Technology (IJERT)},
Jun. 2025.

\bibitem{sivaramakrishnan2025mlcalmo}
R.~Sivaramakrishnan et al.,
``Machine learning-enhanced last-mile delivery optimisation:
integrating deep reinforcement learning with queueing theory for
dynamic vehicle routing,''
\textit{Applied Sciences}, vol.~15, no.~21, p.~11320, Oct. 2025,
doi: 10.3390/app15211320.

%% ── WebSocket Protocols ──────────────────────────────────────────────────

\bibitem{pimentel2012websocket}
V.~Pimentel and B.~G. Nickerson,
``Communicating and displaying real-time data with WebSocket,''
\textit{IEEE Internet Computing}, vol.~16, no.~4, pp.~45--53,
Jul.--Aug. 2012,
doi: 10.1109/MIC.2012.64.

\bibitem{enache2023mqtt}
B.~A. Enache, C.~K. Banica, and A.~G. Bogdan,
``Performance analysis of MQTT over WebSocket for IoT applications,''
\textit{Scientific Bulletin of Electrical Engineering Faculty},
vol.~23, no.~1, pp.~46--49, 2023.

\bibitem{bedir2025websocket}
M.~Hlayel et al.,
``Latency analysis of WebSocket and industrial protocols in IIoT
networks,''
\textit{Int. J. Engineering Trends and Technology (IJETT)},
vol.~73, no.~1, pp.~120--135, Jan. 2025.

%% ── Supporting Reference ─────────────────────────────────────────────────

\bibitem{hughes2015geomesa}
T.~Hughes, D.~Fox, M.~Butler, and P.~Blackshaw,
``GeoMesa: A distributed architecture for spatio-temporal fusion,''
in \textit{Proc. SPIE Defense and Security Symposium}, Baltimore, MD,
USA, 2015, vol.~9473.

\end{thebibliography}

\end{document}